\chapter{Documentação de engenharia de software do PaP}\label{ap:engenharia}
Os diagramas desta revisão documentam retrospectivamente a extensão São Paulo. A notação é simplificada, baseada nas perspectivas de casos de uso, sequência e estados da UML \citep{uml251}. Os nomes dos estados são conceituais; não são uma transcrição de variáveis da implementação. A separação mantém no corpo a visão geral e reserva aqui os detalhes de comportamento.
\section{Atores e casos de uso}
O passageiro utiliza o fluxo de consulta sem cadastro obrigatório. O administrador representa a operação técnica do cadastro; importação por comando não implica existência de um botão na interface do passageiro.

\begin{figure}[H]
\centering\begin{tikzpicture}[font=\small,uc/.style={draw,ellipse,align=center,text width=4.4cm,minimum height=.85cm}]
\node (p) at (0,0) {Passageiro};\node (a) at (0,-6) {Administrador};
\draw (2,1) rectangle (9,-7);\node at (5.5,.65) {PaP};
\foreach \y/\id/\txt in {0/u1/Identificar embarque,-1.2/u2/Consultar próximos ônibus,-2.4/u3/Planejar por destino,-3.6/u4/Ver rota e trechos no mapa,-4.8/u5/Gerenciar favoritos,-6/u6/Importar e consultar cadastro}{\node[uc] (\id) at (5.5,\y-.2) {\txt};}
\foreach \id in {u1,u2,u3,u4,u5}{\draw (p)--(\id.west);}\draw (a)--(u6.west);
\end{tikzpicture}
\caption{Atores e principais casos de uso do PaP.}\label{fig:casos-pap}
\fonte{Elaboração do grupo, documentação da implementação em 7 out. 2026.}
\end{figure}

\section{Consulta progressiva de alternativas}
A sequência representa uma consulta bem-sucedida. O backend pode emitir um primeiro resultado sem caminhada antes de validar alternativas; sua existência não é garantida. A consulta externa pode falhar e excluir conexões ainda não verificadas. A repetição está limitada conforme o Apêndice~\ref{ap:rotas}; seleção de rota e cálculo podem coexistir.

\begin{figure}[H]
\centering\begin{tikzpicture}[>=Stealth,font=\footnotesize]
\foreach \x/\name in {0/Passageiro,3/Interface,6/Backend,9/Roteamento a pé}{\node[draw,align=center,text width=2.3cm] at (\x,0) {\name};\draw[dashed] (\x,-0.4)--(\x,-10);}
\draw[->] (0,-1)--node[above]{Origem e destino}(3,-1);
\draw[->] (3,-2)--node[above]{Solicitar alternativas}(6,-2);
\node[align=center,text width=4cm,fill=white] at (7.6,-3) {Busca por rodadas no GTFS};
\draw[->,dashed] (6,-4)--node[above]{Resultado inicial}(3,-4);
\draw[->] (3,-4.7)--node[above]{Apresentar alternativas}(0,-4.7);
\draw[->] (6,-5.5)--node[above]{Validar caminhada candidata}(9,-5.5);
\draw[->,dashed] (9,-6.3)--node[above]{Distância / falha}(6,-6.3);
\draw[->,dashed] (6,-7.2)--node[above]{Novas alternativas (NDJSON)}(3,-7.2);
\draw[->] (0,-8)--node[above]{Selecionar uma rota}(3,-8);
\node[align=center,text width=3cm] at (3,-9) {Abrir mapa; preservar escolha e zoom nas atualizações};
\draw[draw] (2.6,-5) rectangle (10.4,-7.6);
\node[anchor=south west] at (2.6,-5) {\textbf{loop} [há candidatas e limites disponíveis]};
\end{tikzpicture}
\caption{Sequência simplificada do cálculo e seleção de uma alternativa.}\label{fig:sequencia-pap}
\fonte{Elaboração do grupo, documentação da implementação em 7 out. 2026.}
\end{figure}

\section{Estados da navegação}
Cancelar a confirmação mantém a tela anterior; confirmar retorna à seleção de destino, preservando o embarque e enfocando o destino anterior. A confirmação pode ser aberta tanto pela lista quanto pelo mapa. A visualização e o andamento do cálculo são dimensões distintas. Com localização conhecida, o fluxo inicia na seleção de destino. Alterar o embarque preserva o destino e aceita parada ou coordenada por toque prolongado. Uma falha de atualização não deve ser interpretada como cancelamento pelo usuário. Sem alternativa disponível, a interface permite revisar a consulta; as transições de recuperação complementam o caminho principal ilustrado.

\begin{figure}[H]
\centering\begin{tikzpicture}[>=Stealth,font=\small,box/.style={draw,rounded corners,align=center,text width=3.7cm,minimum height=.85cm,fill=black!3},arr/.style={->,thick},lab/.style={font=\scriptsize,fill=white,align=center,inner sep=2pt}]
\node[circle,fill,minimum size=2.5mm,inner sep=0pt] (start) at (0,1.3) {};
\node[box] (a) at (0,0) {Selecionando embarque};
\node[box] (b) at (6,0) {Consultando próximos ônibus};
\node[box] (c) at (0,-2.6) {Selecionando destino};
\node[box] (e) at (0,-4.4) {Consultando alternativas};
\node[box] (f) at (0,-6.2) {Visualizando rota no mapa};
\node[box] (conf) at (6,-6.2) {Confirmando cancelamento};
\draw[arr] (start)--(a);
\draw[arr] (start.west)--(-3,1.3)--(-3,-2.6)--(c.west);
\node[lab,text width=1.9cm] at (-3,-1.2) {localização\\conhecida};
\draw[arr] (a)--node[lab,above]{QR Code\\ou link}(b);
\draw[arr] (a)--node[lab,right,pos=.3]{embarque\\no mapa}(c);
\draw[arr] (b.south)--(6,-1.7)--node[lab,above]{selecionar destino}(0,-1.7)--(c.north);
\draw[arr] (c)--node[lab,right]{destino definido}(e);
\draw[arr] ([xshift=7pt]e.south)--node[lab,right]{selecionar rota}([xshift=7pt]f.north);
\draw[arr] ([xshift=-7pt]f.north)--node[lab,left]{reabrir lista}([xshift=-7pt]e.south);
\draw[arr] ([yshift=6pt]f.east)--node[lab,above]{cancelar}([yshift=6pt]conf.west);
\draw[arr] ([yshift=-6pt]conf.west)--node[lab,below]{continuar}([yshift=-6pt]f.east);
\draw[arr] (conf.east)--(8.3,-6.2)--(8.3,2)--(1,2)--([xshift=1cm]a.north);
\node[lab,text width=2.3cm] at (7.8,-3.3) {confirmar:\\desmarcar embarque\\e destino};
\draw[arr] (c.west)--(-4.4,-2.6)--(-4.4,0)--(a.west);
\node[lab,text width=1.1cm] at (-4.4,-.9) {sem\\origem};
\draw[arr] (f.south)--(0,-7.3)--(-5.2,-7.3)--(-5.2,.7)--(-1,.7)--([xshift=-1cm]a.north);
\node[lab,text width=2.3cm] at (-5.2,-4.7) {alterar\\embarque:\\preservar destino};
\node[align=center,text width=13cm,font=\scriptsize] at (1.2,-8.5) {A confirmação também pode ser aberta na lista. Continuar mantém a visualização anterior;\\confirmar preserva os favoritos e o enquadramento do mapa.};
\end{tikzpicture}

\caption{Estados conceituais da navegação e confirmação de cancelamento.}\label{fig:estados-pap}
\fonte{Elaboração do grupo, documentação da implementação em 7 out. 2026.}
\end{figure}

\section{Rastreabilidade e evidências}
A Tabela~\ref{tab:rastreabilidade} relaciona requisitos, responsabilidades e verificações localizadas no código. A existência de um teste não significa que ele foi executado nesta revisão documental. Os resultados históricos de execução mantêm as datas e versões já registradas no relatório.
\begin{longtable}{p{1.2cm}p{3cm}p{5cm}p{3cm}}
\caption{Rastreabilidade dos requisitos e limites das evidências.}\label{tab:rastreabilidade}\\
\hline ID & Responsabilidade & Verificação identificada / proposta & Limite\\\hline\endfirsthead
\hline ID & Responsabilidade & Verificação identificada / proposta & Limite\\\hline\endhead
RF01 & Entrada e seleção do ponto & Testes de embarque; completar roteiro de leitura de QR Code em dispositivo físico & Não prova implantação no ponto\\
RF02 & Consulta de partidas & Testes backend em \nolinkurl{test_sptrans.py}; inspecionar aviso da fonte & Disponibilidade real varia\\
RF03 & Planejador e caminhada & \nolinkurl{test_sptrans_planner.py} e \nolinkurl{test_sptrans_walking.py} & Cenários controlados, não validação de toda a rede\\
RF04 & Mapa e lista & Testes E2E de seleção de trecho, baldeações B/C e caminhada final & Compreensão depende de avaliação\\
RF05 & Favoritos locais & Testes E2E de nome persistido, edição e confirmação de remoção & Armazenamento depende do navegador\\
RF06 & Diálogo de cancelamento & Teste E2E de cancelamento confirmado & Não mede erros reais de uso\\
RQ01 & Estado da interface e resposta progressiva & Testes E2E de stream, preservação da escolha e atualização sem redefinir mapa & Não mede tempo em toda condição de rede\\
RQ02 & Fluxo de interação & Protocolo por tarefas descrito na metodologia & Avaliação com passageiros pendente\\\hline
\end{longtable}
Os testes E2E citados estão em \nolinkurl{frontend/e2e/sptrans.spec.js}; os testes backend estão em \nolinkurl{backend/transit/}. Capturas e inspeção de código complementam a documentação, mas não substituem execução de testes, avaliação de acessibilidade ou participação da comunidade.

Na revisão ab32235, a verificação de regressão executou 49 testes de planejamento e caminhada no backend, os quatro arquivos de testes unitários do frontend e um cenário automatizado em emulação móvel. Este último verifica toque prolongado, cancelamento por arrasto e dois dedos e preservação do destino ao alterar o embarque. Na revisão b70873a, os 21 testes unitários do frontend passaram, incluindo três casos de agrupamento de frequências: faixas consecutivas, lacunas ou mudanças de frequência e passagem pela meia-noite. A compilação da aplicação também foi concluída. Esses resultados documentam verificações técnicas, sem equivaler a teste físico no celular ou avaliação de usabilidade com passageiros.
